% !TEX program = xelatex
\documentclass[a4paper,14pt,oneside,openany]{memoir}

%%% Задаем поля, отступы и межстрочный интервал %%%

\usepackage[left=30mm, right=15mm, top=20mm, bottom=20mm]{geometry}
\pagestyle{plain}
\parindent=1.25cm
\usepackage{indentfirst}

%%% Задаем языковые параметры и шрифт %%%

\usepackage[english,russian]{babel}
\babelfont[russian]{rm}{Times New Roman}
\babelfont[russian]{tt}{Courier New}

%%% Задаем стиль заголовков и подзаголовков в тексте %%%

\setsecnumdepth{subsection}
\renewcommand*{\chapterheadstart}{}
\renewcommand*{\printchaptername}{}
\renewcommand*{\chapnumfont}{\normalfont\bfseries}
\renewcommand*{\afterchapternum}{\hspace{1em}}
\renewcommand*{\printchaptertitle}{\normalfont\bfseries\centering\MakeUppercase}
\setbeforesecskip{20pt}
\setaftersecskip{20pt}
\setsecheadstyle{\raggedright\normalfont\bfseries}
\setbeforesubsecskip{20pt}
\setaftersubsecskip{20pt}
\setsubsecheadstyle{\raggedright\normalfont\bfseries}

%%% Отключаем нумерацию секций %%%
\setsecnumdepth{none}

%%% Сквозная нумерация рисунков и таблиц %%%

\counterwithout{figure}{chapter}
\counterwithout{table}{chapter}

%%% Выравнивание и переносы %%%

\tolerance 1414
\hbadness 1414
\emergencystretch 1.5em
\clubpenalty=10000
\widowpenalty=10000

%%% Задаем параметры оформления рисунков и таблиц %%%

\usepackage{graphicx}
\usepackage{caption}
\usepackage{subcaption}
\usepackage{float}
\graphicspath{{photo/}}
\captionsetup[figure]{font=small, width=\textwidth, name=Рисунок, justification=centering}
\captiondelim{ --- }
\usepackage[section]{placeins}

%%% Задаем параметры ссылок и гиперссылок %%% 

\usepackage{hyperref}
\hypersetup{
    colorlinks=true,
    linkcolor=blue,
    citecolor=blue,
    urlcolor=blue
}

%%% Настраиваем отображение списков %%%

\usepackage{enumitem}
\renewcommand*{\labelitemi}{\normalfont{--}}
\setlist{noitemsep, leftmargin=*}

\begin{document}

% Титульный лист (без номера страницы)
\thispagestyle{empty}

\begin{center}
    МИНИСТЕРСТВО ЦИФРОВОГО РАЗВИТИЯ, СВЯЗИ И МАССОВЫХ \\ КОММУНИКАЦИЙ \\ РОССИЙСКОЙ ФЕДЕРАЦИИ

    \vspace{20pt}

    Федеральное государственное бюджетное образовательное учреждение \\ высшего образования \\
    «Сибирский государственный университет телекоммуникаций и \\ информатики»

    \vspace{20pt}

    Кафедра телекоммуникационных систем и вычислительных средств \\ (ТСиВС)
\end{center}

\vfill

\begin{center}
    Лабораторная работа №13 \\
    по дисциплине \\
    \textit{«WEB-Технологии»}

    \vspace{20pt}

    по теме: \\
    \uppercase{Погружение в разработку на JavaScript}
\end{center}

\vfill

\noindent Студент: \\
\textit{Группа № ИКС-532 \hfill Лёшин А.С.}

\vspace{20pt}

\noindent Преподаватель: \\
\textit{старший преподаватель \hfill Андреев А.В.}

\vfill

\begin{center}
    Новосибирск 2026 г.
\end{center}

\newpage
\setcounter{page}{2}
\OnehalfSpacing*

\section{Блок 1 Введение}

Целью лабораторной работы является изучение базовых концепций языка JavaScript, включая работу с переменными, массивами, циклами и методами строк. Также в работу входит освоение навыков взаимодействия с внешними API и деплоя веб-приложений на хостинг.

\section{Блок 2 Описание работы}

\section{Блок 2.1 Базовые понятия и бонусная система}

В рамках первой задачи была реализована имитация бонусной системы интернет-магазина. Были созданы переменные для имени пользователя и баланса, а также произведен расчет остатка баллов через 7 дней с учетом ежедневного сгорания и периодических покупок.

\section{Блок 2.2 Основы JavaScript (Массивы и циклы)}

Для имитации мессенджера был инициализирован массив сообщений. С помощью цикла for реализован вывод переписки с чередованием отправителей («Друг»/«Вы»). Также добавлена функция поиска по ключевому слову с использованием метода \texttt{.includes()}.

\begin{figure}[H]
    \centering
    \includegraphics[width=0.7\textwidth]{1.png}
    \caption{Реализация логики бонусной системы и чата}
\end{figure}

\begin{figure}[H]
    \centering
    \includegraphics[width=0.7\textwidth]{2.png}
    \caption{Вывод программы}
\end{figure}

\section{Блок 3 Разработка сайта и работа с API}

На данном этапе было реализовано динамическое получение данных о персонажах Marvel с внешнего API. Полученные данные были преобразованы в интерактивные карточки персонажей.

\begin{figure}[H]
    \centering
    \includegraphics[width=0.7\textwidth]{3.png}
    \caption{Отображение персонажей Marvel, полученных через API}
\end{figure}

\section{Блок 4 Доработка и публикация}

Итоговый проект был загружен на хостинг Timeweb. Файлы проекта (index.html, start.js, index.js) были размещены в директории public\_html.

\begin{figure}[H]
    \centering
    \includegraphics[width=0.7\textwidth]{4.png}
    \caption{Вывод рабочего сайта, загруженного через Timeweb}
\end{figure}

\section{Выводы}

В ходе лабораторной работы были освоены фундаментальные принципы JavaScript: работа с типами данных, управляющими конструкциями и массивами. Были получены практические навыки работы с асинхронными запросами (Fetch API) и размещения статических сайтов на облачном хостинге.

\section{Ссылки на ресурсы}

\begin{itemize}
    \item \textbf{Исходный код (Replit):} \href{https://replit.com/join/fyandlefiy-lesinrtm}{https://replit.com/join/fyandlefiy-lesinrtm}
    \item \textbf{Исходный код с Marvel API (Replit):} \href{https://lab-13--lesinrtm.replit.app}{https://lab-13--lesinrtm.replit.app}
    \item \textbf{Рабочий сайт (Timeweb):} \href{https://ci977889.tw1.ru/}{https://ci977889.tw1.ru/}
\end{itemize}

\end{document}